\section{Discussion}\label{sec:discussion}
The four recognition models achieved similar overall scores, leaving no clear architectural winner for local field recognition. Learned grouping improved the assignment of the same human-annotated entities compared with fixed rules. For Magda, this highlights the importance of both stages: recognising a price and linking it to the right product. In the earlier feature experiments, combining additional spatial context and lexical cues gave the best grouping result against automatic annotations. Adding anchor features reduced the score. More input features did not automatically make the model better.

The complete pipeline addresses our original question more directly. For product names and regular prices, Magda ran locally and reached F1 scores close to Qwen (Table~\ref{tab:offers}). Qwen3.8 recovered more reference records, while Magda had fewer unmatched outputs relative to its total output. Some additional Qwen records were duplicates, while others appeared to be valid offers or variants. Astra scored highest in the separate Codex run, which used shared context and tools. The paired intervals against Magda leave the ranking uncertain, and similar scores do not prove equivalent performance. A record can also fail matching because one field differs, even when other details are correct. These scores therefore give only part of the picture of practical extraction quality.

The project supports the idea that local extraction is viable for product names and regular prices on the evaluated PENNY pages. Magda reached 0.7571 F1, close to Qwen3.8's 0.7643 (Table~\ref{tab:offers}). Processing prepared pages took 0.695 seconds per page on average, versus 26.965 seconds through the Qwen3.8 API (Section~\ref{sec:record-results}). These timings reflect different hardware and network conditions. Annotation, training, maintenance and inference still require resources. The comparison does not show that the pipeline is free or cheaper than using an API overall.
